\chapter{Análisis exploratorio de datos}
\label{ch:eda}

Este capítulo corresponde al notebook \texttt{01\_eda}. Se describe la
distribución de clases, la composición por plataforma y la longitud de los
textos tokenizados.

\section{Distribución de clases}

La Figura~\ref{fig:dist-clases} muestra el conteo por clase en el corpus.
Existe desbalance: la clase Lenguaje Ofensivo es la más frecuente, mientras que
Amenazas/Violencia es minoritaria. Este desbalance motiva el uso de pesos de
clase y la elección de F1-macro como métrica principal.

\begin{figure}[H]
\centering
\includegraphics[width=0.85\textwidth]{01_distribucion_clases.png}
\caption{Distribución de clases en el corpus.}
\label{fig:dist-clases}
\end{figure}

\section{Baselines triviales}

Antes de entrenar modelos, se estiman baselines triviales sobre el
\textbf{conjunto de test} ($n=460$; Tabla~\ref{tab:trivial}): clase mayoritaria,
azar uniforme y azar estratificado. Para la clase mayoritaria se simula un
clasificador que siempre predice la etiqueta más frecuente: precisión $p$
(proporción en test), recall $1{,}0$; el F1 de las demás clases es $0$. El
F1-macro es la media sobre $k=4$ clases. Los valores del CSV se redondean a
tres decimales (\texttt{round}); la tabla generada muestra cuatro. El informe
de preparación v1.1 reportaba $\approx 0{,}129$ por un redondeo intermedio de
$p$; aquí $0{,}128$. El azar se sitúa en $0{,}25$. Cualquier modelo útil debe
superar claramente estos umbrales.

\begin{table}[H]
\centering
\caption{Baselines triviales (F1-macro).}
\label{tab:trivial}
\begin{tabular}{lc}
\toprule
Baseline & F1-macro \\
\midrule
azar\_estratificado & 0.2500 \\
azar\_uniforme & 0.2500 \\
clase\_mayoritaria & 0.1284 \\
\bottomrule
\end{tabular}

\end{table}

\section{Plataforma y longitud de tokens}

La Figura~\ref{fig:plataforma} desglosa las clases por plataforma. X concentra
más volumen; Facebook presenta una distribución distinta que, como se verá en
la evaluación por cortes, resulta más difícil para el modelo.

\begin{figure}[H]
\centering
\includegraphics[width=0.85\textwidth]{02_plataforma_clases.png}
\caption{Distribución de clases por plataforma.}
\label{fig:plataforma}
\end{figure}

La Figura~\ref{fig:tokens} muestra la longitud en tokens del texto de entrada.
La mayoría de instancias cabe cómodamente bajo el \texttt{max\_length} de 128
tokens usado en mBERT.

\begin{figure}[H]
\centering
\includegraphics[width=0.75\textwidth]{03_token_lengths.png}
\caption{Distribución de longitudes en tokens.}
\label{fig:tokens}
\end{figure}
